\documentclass[aspectratio=169,11pt]{beamer}
\usetheme{default}
\usefonttheme{professionalfonts}
\usepackage[T1]{fontenc}
\usepackage[utf8]{inputenc}
\usepackage{lmodern}
\usepackage{amsmath,amssymb,mathtools}
\usepackage{booktabs,array,multicol}
\usepackage{tikz}
\usetikzlibrary{arrows.meta,positioning,calc,shapes.geometric}
\usepackage{listings}
\usepackage{xcolor}
\usepackage{graphicx}
\usepackage{pifont}

\definecolor{pidkblue}{HTML}{0F4165}
\definecolor{pidkblue2}{HTML}{173F5F}
\definecolor{pidklight}{HTML}{EAF3FA}
\definecolor{pidklighter}{HTML}{F5F9FC}
\definecolor{pidkteal}{HTML}{0EA5A8}
\definecolor{pidkgray}{HTML}{2F3B44}
\definecolor{pidkgold}{HTML}{B88700}
\definecolor{codebg}{HTML}{F7F9FB}
\definecolor{codeframe}{HTML}{C9D3DD}

\setbeamertemplate{navigation symbols}{}
\setbeamercolor{normal text}{fg=pidkgray,bg=white}
\setbeamercolor{frametitle}{fg=white,bg=pidkblue}
\setbeamercolor{title}{fg=pidkblue}
\setbeamertemplate{frametitle}{%
  \nointerlineskip
  \begin{beamercolorbox}[wd=\paperwidth,ht=0.86cm,dp=0.27cm,leftskip=0.35cm]{frametitle}%
    \usebeamerfont{frametitle}\insertframetitle
  \end{beamercolorbox}%
}

\newcommand{\coursefooter}{Pemrograman untuk Ilmu Data dan Komputasi}
\newcommand{\programfooter}{Program Studi Matematika ITERA}
\setbeamertemplate{footline}{%
\ifnum\value{framenumber}>1
  \leavevmode%
  \hbox{%
  \begin{beamercolorbox}[wd=.38\paperwidth,ht=2.6ex,dp=1ex,leftskip=0.25cm]{author in head/foot}%
    \scriptsize\color{white}\coursefooter
  \end{beamercolorbox}%
  \begin{beamercolorbox}[wd=.44\paperwidth,ht=2.6ex,dp=1ex,center]{title in head/foot}%
    \scriptsize\color{white}\programfooter
  \end{beamercolorbox}%
  \begin{beamercolorbox}[wd=.18\paperwidth,ht=2.6ex,dp=1ex,rightskip=0.25cm plus1fil]{date in head/foot}%
    \scriptsize\color{white}\hfill\insertframenumber/\inserttotalframenumber
  \end{beamercolorbox}}%
  \vskip0pt%
\fi}
\setbeamercolor{author in head/foot}{bg=pidkblue}
\setbeamercolor{title in head/foot}{bg=pidkblue}
\setbeamercolor{date in head/foot}{bg=pidkblue}

\lstdefinestyle{py}{
  language=Python,
  basicstyle=\ttfamily\scriptsize,
  keywordstyle=\color{pidkblue}\bfseries,
  stringstyle=\color{teal!70!black},
  commentstyle=\color{gray!75!black},
  showstringspaces=false,
  breaklines=true,
  columns=fullflexible,
  keepspaces=true,
  backgroundcolor=\color{codebg},
  frame=single,
  rulecolor=\color{codeframe},
  framerule=0.4pt,
  xleftmargin=0.5mm,
  xrightmargin=0.5mm,
  aboveskip=0.5em,
  belowskip=0.3em
}
\lstset{style=py}

\setbeamertemplate{blocks}[rounded][shadow=true]
\newenvironment{card}[1]{\begin{block}{#1}}{\end{block}}
\setbeamercolor{block title}{fg=white,bg=pidkblue}
\setbeamercolor{block body}{fg=pidkgray,bg=pidklight}
\newenvironment{goldcard}[1]{\begingroup\setbeamercolor{block title}{fg=white,bg=pidkgold}\setbeamercolor{block body}{fg=pidkgray,bg=yellow!7}\begin{block}{#1}}{\end{block}\endgroup}
\newcommand{\smallnote}[1]{\vspace{0.10cm}\begin{beamercolorbox}[wd=\linewidth,sep=0.14cm]{note}\small #1\end{beamercolorbox}}
\setbeamercolor{note}{fg=pidkgray,bg=pidklighter}

\title{Pertemuan 9: Regresi Linier, Least Squares, dan Pandas}
\author{\textbf{Triyana Muliawati, S.Si., M.Si.}\\Rifky Fauzi}
\date{Semester Ganjil 2026/2027}

\begin{document}

% 1
{\setbeamertemplate{footline}{}
\begin{frame}[plain]
\begin{tikzpicture}[remember picture,overlay]
  \fill[pidkgray!96] (current page.south west) rectangle ([xshift=0.48\paperwidth]current page.north west);
  \fill[white] ([xshift=0.48\paperwidth]current page.south west) rectangle (current page.north east);
  \draw[pidkteal,line width=1.3pt] ([xshift=0.05\paperwidth,yshift=0.24\paperheight]current page.south west) -- ([xshift=0.42\paperwidth,yshift=0.64\paperheight]current page.south west);
  \foreach \x/\y in {0.07/0.31,0.12/0.34,0.18/0.42,0.23/0.39,0.29/0.52,0.35/0.54,0.40/0.63}{
    \fill[white] ([xshift=\x\paperwidth,yshift=\y\paperheight]current page.south west) circle (2pt);
  }
  \node[anchor=west,align=left,text=pidkblue] at ([xshift=0.53\paperwidth,yshift=0.80\paperheight]current page.south west) {\parbox{0.42\paperwidth}{\fontsize{28}{31}\selectfont Pertemuan 9}};
  \node[anchor=west,align=left,text=pidkblue] at ([xshift=0.53\paperwidth,yshift=0.63\paperheight]current page.south west) {\parbox{0.42\paperwidth}{\fontsize{20}{23}\selectfont Regresi Linier,\\Least Squares,\\dan Pandas}};
  \node[anchor=west,align=left,text=black] at ([xshift=0.53\paperwidth,yshift=0.40\paperheight]current page.south west) {\parbox{0.42\paperwidth}{\fontsize{12}{15}\selectfont Pemrograman untuk Ilmu Data\\dan Komputasi\\[0.08cm]MA25-21008}};
  \node[anchor=west,align=left,text=black] at ([xshift=0.53\paperwidth,yshift=0.16\paperheight]current page.south west) {\parbox{0.42\paperwidth}{\fontsize{11}{14}\selectfont \textbf{Triyana Muliawati, S.Si., M.Si.}\\Rifky Fauzi\\[0.16cm]Program Studi Matematika\\Institut Teknologi Sumatera\\[0.12cm]Semester Ganjil 2026/2027}};
\end{tikzpicture}
\end{frame}}

% 2
\begin{frame}{Era Big Data dan Peran Matematika}
\small
\begin{columns}[T,totalwidth=\textwidth]
\column{0.48\textwidth}
\begin{card}{Data sebagai objek komputasi}
Data modern dapat berukuran besar, banyak variabel, cepat bertambah, dan berasal dari berbagai sumber:
\begin{itemize}\itemsep0.28em
  \item transaksi dan aktivitas digital;
  \item sensor, cuaca, dan pengamatan lingkungan;
  \item proses industri dan produksi;
  \item data pendidikan, ekonomi, dan sosial.
\end{itemize}
\end{card}
\column{0.49\textwidth}
\begin{goldcard}{Peran matematika}
Matematika memberi kerangka untuk
\begin{itemize}\itemsep0.28em
  \item merepresentasikan data sebagai vektor dan matriks;
  \item merumuskan hubungan antarvariabel;
  \item menentukan ukuran galat atau fungsi objektif;
  \item mencari parameter melalui optimisasi;
  \item memeriksa kewajaran hasil komputasi.
\end{itemize}
\end{goldcard}
\end{columns}
\vspace{0.15cm}
\centering
\textbf{Data $\longrightarrow$ representasi matematis $\longrightarrow$ algoritma $\longrightarrow$ informasi}
\end{frame}

% 3
\begin{frame}{Pemodelan Hubungan Antarvariabel}
\small
\begin{columns}[T,totalwidth=\textwidth]
\column{0.48\textwidth}
\begin{card}{Pertanyaan matematis}
Untuk pasangan pengamatan $(x_i,y_i)$, matematika dapat digunakan untuk menjawab:
\begin{itemize}\itemsep0.30em
  \item apakah perubahan $x$ berkaitan dengan perubahan $y$?
  \item seberapa besar perubahan $y$ ketika $x$ berubah?
  \item bagaimana memperkirakan $y$ untuk $x$ yang belum diamati?
\end{itemize}
Model umum:
\[
 y_i=f(x_i)+\varepsilon_i.
\]
\end{card}
\column{0.49\textwidth}
\begin{goldcard}{Model linier sebagai pendekatan awal}
Pilih fungsi paling sederhana
\[
 f(x)=b_0+b_1x,
\]
sehingga
\[
 y_i=b_0+b_1x_i+\varepsilon_i.
\]
Menentukan $b_0$ dan $b_1$ dari data disebut \textbf{regresi linier sederhana}.
\end{goldcard}
\vspace{0.12cm}
\textbf{Contoh:} jam belajar $\to$ nilai; luas bangunan $\to$ harga; suhu $\to$ konsumsi listrik; curah hujan $\to$ tinggi muka air.
\end{columns}
\end{frame}

% 4
\begin{frame}{Regresi Linier Sederhana dan Metode Kuadrat Terkecil}
\begin{columns}[T,totalwidth=\textwidth]
\column{0.47\textwidth}
\begin{card}{Residual dan fungsi objektif}
Prediksi untuk titik ke-$i$:
\[
\widehat y_i=b_0+b_1x_i.
\]
Residual:
\[
e_i=y_i-\widehat y_i.
\]
Metode kuadrat terkecil memilih $b_0,b_1$ yang meminimumkan
\[
S(b_0,b_1)=\sum_{i=1}^{n}(y_i-b_0-b_1x_i)^2.
\]
\end{card}
\column{0.50\textwidth}
\begin{goldcard}{Kondisi minimum}
Dengan
\[
\frac{\partial S}{\partial b_0}=0,\qquad
\frac{\partial S}{\partial b_1}=0,
\]
diperoleh
\[
 b_1=\frac{\sum (x_i-\bar x)(y_i-\bar y)}{\sum (x_i-\bar x)^2},
\qquad
 b_0=\bar y-b_1\bar x.
\]
\end{goldcard}
\smallnote{Rumus koefisien bukan rumus yang muncul begitu saja: ia merupakan solusi masalah minimisasi jumlah kuadrat residual.}
\end{columns}
\end{frame}

% 5
\begin{frame}[fragile]{Algoritma Regresi Linier Sederhana}
\begin{columns}[T,totalwidth=\textwidth]
\column{0.46\textwidth}
\begin{card}{Prosedur komputasi}
\begin{enumerate}\itemsep0.36em
  \item Hitung $\bar x$ dan $\bar y$.
  \item Hitung $N=\sum (x_i-\bar x)(y_i-\bar y)$.
  \item Hitung $D=\sum (x_i-\bar x)^2$.
  \item Jika $D=0$, model tidak dapat menentukan kemiringan.
  \item Hitung $b_1=N/D$ dan $b_0=\bar y-b_1\bar x$.
  \item Gunakan $\widehat y=b_0+b_1x$ untuk prediksi.
\end{enumerate}
\end{card}
\column{0.51\textwidth}
\begin{lstlisting}
def fit_linear(x, y):
    n = len(x)
    xbar = sum(x) / n
    ybar = sum(y) / n

    num, den = 0.0, 0.0
    for xi, yi in zip(x, y):
        num += (xi-xbar)*(yi-ybar)
        den += (xi-xbar)**2

    if den == 0:
        raise ValueError("x tidak bervariasi")

    b1 = num / den
    b0 = ybar - b1*xbar
    return b0, b1
\end{lstlisting}
\end{columns}
\end{frame}

% 6
\begin{frame}[fragile]{Representasi dan Eksplorasi Data Tabular}
\begin{columns}[T,totalwidth=\textwidth]
\column{0.50\textwidth}
\begin{lstlisting}
import pandas as pd

df = pd.read_csv("data_regresi_m09.csv")

print(df.head())
print(df.shape)
print(df.dtypes)
print(df.isna().sum())

# visualisasi hubungan satu prediktor
df.plot.scatter(
    x="jam_belajar",
    y="nilai",
    grid=True
)
\end{lstlisting}
\column{0.47\textwidth}
\begin{card}{Struktur data}
Dataset latihan mempunyai tiga kolom:
\begin{itemize}\itemsep0.38em
  \item \texttt{jam\_belajar}: prediktor pertama;
  \item \texttt{jumlah\_latihan}: prediktor kedua;
  \item \texttt{nilai}: variabel respons.
\end{itemize}
Pada tahap awal, gunakan hanya \texttt{jam\_belajar} untuk regresi sederhana.
\end{card}
\begin{goldcard}{Tujuan eksplorasi}
Periksa tipe data, nilai hilang, rentang nilai, dan pola visual sebelum melakukan estimasi parameter.
\end{goldcard}
\end{columns}
\end{frame}

% 7
\begin{frame}[fragile]{Estimasi Parameter dan Evaluasi Model}
\small
\begin{columns}[T,totalwidth=\textwidth]
\column{0.47\textwidth}
\begin{card}{Tujuan pemisahan data}
Model dibangun agar pola yang dipelajari dapat digunakan pada data yang tidak ikut menentukan parameter.
\begin{itemize}\itemsep0.28em
  \item \textbf{Training set}: mengestimasi parameter model.
  \item \textbf{Test set}: menilai kinerja model setelah parameter ditentukan.
\end{itemize}
Salah satu ukuran galat pada test set:
\[
\mathrm{MSE}=\frac{1}{m}\sum_{i=1}^{m}(y_i-\widehat y_i)^2.
\]
\end{card}
\column{0.50\textwidth}
\begin{lstlisting}[basicstyle=\ttfamily\tiny]
import numpy as np

idx = np.arange(len(df))
rng = np.random.default_rng(7)
rng.shuffle(idx)

cut = int(0.7 * len(df))
train = df.iloc[idx[:cut]]
test  = df.iloc[idx[cut:]]

b0, b1 = fit_linear(
    train["jam_belajar"].to_numpy(),
    train["nilai"].to_numpy()
)

yhat = b0 + b1*test["jam_belajar"].to_numpy()
mse = np.mean((test["nilai"].to_numpy()-yhat)**2)
\end{lstlisting}
\end{columns}
\end{frame}

% 8
\begin{frame}{Regresi Linier Berganda}
\begin{columns}[T,totalwidth=\textwidth]
\column{0.47\textwidth}
\begin{card}{Dua prediktor}
Misalkan nilai dipengaruhi oleh
\[
x_1=\text{jam belajar},
\]
\[
x_2=\text{jumlah latihan}.
\]
Model liniernya
\[
\widehat y=\beta_0+\beta_1x_1+\beta_2x_2.
\]
Setiap koefisien menyatakan kontribusi satu prediktor ketika prediktor lain dipertahankan.
\end{card}
\column{0.50\textwidth}
\begin{goldcard}{Dari skalar ke matriks}
Untuk $n$ observasi, susun
\[
X=\begin{bmatrix}
1 & x_{11} & x_{12}\\
1 & x_{21} & x_{22}\\
\vdots & \vdots & \vdots\\
1 & x_{n1} & x_{n2}
\end{bmatrix},\qquad
\boldsymbol\beta=\begin{bmatrix}\beta_0\\\beta_1\\\beta_2\end{bmatrix}.
\]
Maka semua prediksi sekaligus ditulis
\[
\widehat{\mathbf y}=X\boldsymbol\beta.
\]
\end{goldcard}
\end{columns}
\end{frame}

% 9
\begin{frame}{Formulasi Least Squares dalam Bentuk Matriks}
\scriptsize
\begin{columns}[T,totalwidth=\textwidth]
\column{0.51\textwidth}
\begin{card}{Dimensi dan norma Euclid}
Untuk dua prediktor:
\[
X\in\mathbb{R}^{n\times 3},\quad
\boldsymbol\beta\in\mathbb{R}^{3\times 1},\quad
\mathbf y\in\mathbb{R}^{n\times 1}.
\]
Residual:
\[
\mathbf e=\mathbf y-X\boldsymbol\beta\in\mathbb{R}^{n\times 1}.
\]
Norma-2 adalah panjang Euclid vektor:
\[
\|\mathbf e\|_2=\sqrt{e_1^2+\cdots+e_n^2}.
\]
Maka
\[
\|\mathbf y-X\boldsymbol\beta\|_2^2=\sum_{i=1}^{n}e_i^2.
\]
\end{card}
\column{0.46\textwidth}
\begin{goldcard}{Masalah least squares}
\[
\min_{\boldsymbol\beta}\;\|\mathbf y-X\boldsymbol\beta\|_2^2.
\]
Kondisi minimum menghasilkan
\[
X^TX\boldsymbol\beta=X^T\mathbf y.
\]
Jika $X^TX$ nonsingular,
\[
\widehat{\boldsymbol\beta}=(X^TX)^{-1}X^T\mathbf y.
\]
\end{goldcard}
\begin{card}{Pertanyaan dimensi}
Mengapa $\widehat{\mathbf y}=X\boldsymbol\beta$, bukan $\boldsymbol\beta X$?\\[0.05cm]
\textit{Hint:} periksa ukuran pada aturan $(a\times b)(b\times c)$.
\end{card}
\end{columns}
\end{frame}

% 10
\begin{frame}[fragile]{Implementasi Regresi Linier Berganda}
\begin{columns}[T,totalwidth=\textwidth]
\column{0.53\textwidth}
\begin{lstlisting}
import numpy as np

X0 = df[["jam_belajar",
         "jumlah_latihan"]].to_numpy(float)
y = df["nilai"].to_numpy(float)

n = len(df)
X = np.column_stack([np.ones(n), X0])

A = X.T @ X
c = X.T @ y
beta = np.linalg.solve(A, c)

yhat = X @ beta
res = y - yhat
mse = np.mean(res**2)
\end{lstlisting}
\column{0.44\textwidth}
\begin{card}{Urutan algoritma}
\begin{enumerate}\itemsep0.34em
  \item Bentuk matriks desain $X$.
  \item Hitung $A=X^TX$.
  \item Hitung $c=X^T\mathbf y$.
  \item Selesaikan $A\boldsymbol\beta=c$.
  \item Hitung $\widehat{\mathbf y}=X\boldsymbol\beta$.
  \item Hitung residual dan galat.
\end{enumerate}
\end{card}
\begin{goldcard}{Hasil pada data contoh}
Dengan seluruh 10 observasi,
\[
\widehat y\approx45.18+3.99x_1+1.97x_2.
\]
\end{goldcard}
\end{columns}
\end{frame}

% 11
\begin{frame}{Model Machine Learning dan Prinsip Optimisasi}
\small
\begin{columns}[T,totalwidth=\textwidth]
\column{0.48\textwidth}
\begin{card}{Prinsip umum pembelajaran model}
Banyak metode machine learning dapat dipandang sebagai proses:
\[
\text{pilih }f(x;\theta)
\to \text{ukur kesalahan}
\to \text{tentukan }\theta.
\]
Pada regresi linier,
\[
L(\beta)=\|y-X\beta\|_2^2.
\]
Pada model lain, bentuk model dan fungsi objektif dapat berbeda.
\end{card}
\column{0.49\textwidth}
\begin{goldcard}{Beberapa keluarga model}
\begin{itemize}\itemsep0.35em
  \item \textbf{Regresi:} linear regression, polynomial regression, regression tree, neural network.
  \item \textbf{Klasifikasi:} logistic regression, k-NN, decision tree, SVM, neural network.
  \item \textbf{Pengelompokan:} k-means, hierarchical clustering.
  \item \textbf{Reduksi dimensi:} PCA.
\end{itemize}
\end{goldcard}
\end{columns}
\smallnote{Fokus mata kuliah: memahami representasi matematis, prosedur komputasi, dan evaluasi model sebelum menggunakan library pemodelan.}
\end{frame}

% 12
\begin{frame}{Pengayaan: Supervised Learning, Prediction, dan Forecasting}
\small
\begin{columns}[T,totalwidth=\textwidth]
\column{0.47\textwidth}
\begin{card}{Supervised learning}
Data tersedia sebagai pasangan
\[
(\mathbf x_i,y_i).
\]
Target $y$ diketahui pada data pelatihan dan dipakai untuk mempelajari hubungan dari prediktor $\mathbf x$.
\begin{itemize}\itemsep0.22em
  \item $y$ numerik $\to$ regresi;
  \item $y$ kategori $\to$ klasifikasi.
\end{itemize}
\end{card}
\column{0.50\textwidth}
\begin{goldcard}{Prediction dan forecasting}
\textbf{Prediction} adalah perkiraan target yang belum diketahui berdasarkan model.\\[0.15cm]
\textbf{Forecasting} adalah prediction untuk nilai masa depan pada data yang mempunyai urutan waktu. Karena masa depan tidak boleh digunakan untuk melatih masa lalu, pembagian data dilakukan secara kronologis.
\end{goldcard}
\begin{card}{Hubungan dengan pertemuan ini}
Regresi linier merupakan salah satu model supervised learning untuk target numerik.
\end{card}
\end{columns}
\end{frame}

% 13
\begin{frame}{Latihan: Regresi Linier Berganda dengan Dua Prediktor}
\scriptsize
\begin{columns}[T,totalwidth=\textwidth]
\column{0.46\textwidth}
\begin{card}{Data}
\centering
\begin{tabular}{c|rrr}
\toprule
$i$ & $x_1$ jam & $x_2$ latihan & $y$ nilai\\\midrule
1 & 1 & 2 & 54\\
2 & 2 & 5 & 61\\
3 & 3 & 1 & 61\\
4 & 4 & 6 & 72\\
5 & 5 & 3 & 71\\
6 & 6 & 7 & 84\\
7 & 7 & 4 & 79\\
8 & 8 & 8 & 95\\
9 & 9 & 5 & 90\\
10&10 & 9 &103\\
\bottomrule
\end{tabular}
\end{card}
\column{0.51\textwidth}
\begin{goldcard}{Tugas}
Bangun model
\[
\widehat y=\beta_0+\beta_1x_1+\beta_2x_2
\]
\textbf{tanpa library pemodelan}. Tuliskan algoritma yang Anda gunakan, kemudian:
\begin{enumerate}\itemsep0.24em
  \item bentuk $X$ dan $y$, beserta dimensinya;
  \item hitung $X^TX$ dan $X^Ty$;
  \item selesaikan persamaan normal untuk memperoleh $\beta$;
  \item hitung prediksi dan residual setiap data;
  \item hitung MSE;
  \item tuliskan $T(n)$ dari prosedur Anda dan tentukan Big-O, dengan $n$ = jumlah observasi dan jumlah prediktor tetap dua.
\end{enumerate}
\end{goldcard}
\end{columns}
\end{frame}

\end{document}

